\documentclass[12pt, letterpaper]{article}

% ── Paquetes ─────────────────────────────────────────────────────────────────
\usepackage[utf8]{inputenc}
\usepackage[T1]{fontenc}
\usepackage[spanish]{babel}
\usepackage{geometry}
\usepackage{graphicx}
\usepackage{booktabs}
\usepackage{hyperref}
\usepackage{float}
\usepackage{caption}
\usepackage{subcaption}
\usepackage{enumitem}
\usepackage{parskip}
\usepackage{titlesec}
\usepackage{fancyhdr}
\usepackage{xcolor}
\usepackage{lmodern}
\usepackage{microtype}
\usepackage{csquotes}

% ── Geometría ─────────────────────────────────────────────────────────────────
\geometry{
  top=2.5cm, bottom=2.5cm,
  left=2.5cm, right=2.5cm
}
\setlength{\headheight}{13.6pt}

% ── Hyperref ──────────────────────────────────────────────────────────────────
\hypersetup{
  colorlinks=true,
  linkcolor=black,
  urlcolor=blue!70!black,
  citecolor=black
}

% ── Encabezado / pie de página ────────────────────────────────────────────────
\pagestyle{fancy}
\fancyhf{}
\fancyhead[L]{\small Visualización de Datos -- Semestre I, 2026}
\fancyhead[R]{\small Tarea 3}
\fancyfoot[C]{\thepage}
\renewcommand{\headrulewidth}{0.4pt}

% ── Título de secciones ───────────────────────────────────────────────────────
\titleformat{\section}{\large\bfseries}{}{0em}{\thesection.\quad}
\titleformat{\subsection}{\normalsize\bfseries}{}{0em}{\thesubsection.\quad}

% ─────────────────────────────────────────────────────────────────────────────
\begin{document}

% ── Portada ───────────────────────────────────────────────────────────────────
\begin{titlepage}
  \centering
  \vspace*{1.5cm}

  {\Large\bfseries Universidad Técnica Federico Santa María}\\[0.3cm]
  {\large Departamento de Informática}\\[0.2cm]
  {\large Visualización de Datos -- Semestre I, 2026}

  \vspace{1.5cm}
  \rule{\linewidth}{0.5pt}\\[0.4cm]
  {\LARGE\bfseries Tarea 3: Mapas}\\[0.2cm]
  {\large Macrotema: \textit{Muerte}}
  \rule{\linewidth}{0.5pt}

  \vspace{1.5cm}

  \begin{tabular}{ll}
    \textbf{Integrante 1:} & Alexis Mellis \\[0.2cm]
    \textbf{Integrante 2:} & Andres Aguila \\
  \end{tabular}

  \vspace{1cm}

  \begin{tabular}{ll}
    \textbf{Profesores:}  & José Luis Martí Lara \\[0.2cm]
    \textbf{Ayudantes:}   & Fernanda Araya Zárate \\
                          & Carlos Arévalo Guajardo \\
  \end{tabular}

  \vfill

\end{titlepage}

% ── Índice ────────────────────────────────────────────────────────────────────
\tableofcontents
\newpage

% ─────────────────────────────────────────────────────────────────────────────
\section{Introducción}

El presente informe corresponde a la Tarea~3 del curso de Visualización de Datos, que
continúa el análisis del macrotema \textbf{[Macrotema]}. En esta entrega se desarrollan
visualizaciones basadas en mapas, cada una construida con una herramienta distinta y
de tipo diferente, con el objetivo de responder preguntas concretas sobre los datos
seleccionados. El informe documenta el proceso individual de cada integrante, las
decisiones de diseño adoptadas y las conclusiones derivadas de cada mapa.

% ─────────────────────────────────────────────────────────────────────────────
\section{Trabajo individual -- Integrante 1: [Nombre]}

\subsection{Dataset seleccionado}

% Descripción cualitativa del dataset: origen, variables relevantes, cobertura
% geográfica/temporal, y por qué es adecuado para una visualización de mapa.

\begin{itemize}
  \item \textbf{Nombre del dataset:} [nombre]
  \item \textbf{Fuente:} [URL o referencia]
  \item \textbf{Variables relevantes:} [lista breve]
  \item \textbf{Cobertura:} [países, período temporal, granularidad geográfica]
  \item \textbf{Descripción cualitativa:} [2--3 oraciones describiendo qué contiene
        el dataset y sus características principales]
\end{itemize}

\subsection{Objetivo de la visualización}

\begin{itemize}
  \item \textbf{Pregunta que busca responder:} [¿Qué pregunta responde el mapa?]
  \item \textbf{Público objetivo:} [¿A quién está dirigido?]
  \item \textbf{Acción o decisión que podría apoyar:} [¿Qué decisión informaría?]
\end{itemize}

\subsection{Tipo de mapa y herramienta}

\begin{itemize}
  \item \textbf{Tipo de mapa:} [Coropleta / Puntos / Burbujas / Flujo / Calor / otro]
  \item \textbf{Herramienta:} [Flourish / Carto / Kepler.gl / Datawrapper / otra]
  \item \textbf{Justificación del tipo:} [¿Por qué este tipo de mapa es el más adecuado
        para la pregunta y los datos?]
\end{itemize}

\subsection{Esquemática utilizada}

% Indique la lógica narrativa o estructural del mapa: ¿es descriptiva, secuencial,
% relacional, narrativa, comparativa? ¿Cómo guía la atención del lector?

[Ejemplo: ``El mapa sigue una esquemática descriptiva: muestra la distribución
geográfica de un fenómeno sin establecer una secuencia temporal ni relaciones
causales. El objetivo es permitir al lector identificar patrones espaciales de
forma exploratoria...'']

\subsection{Colores utilizados y justificación}

% Mencione la paleta concreta y explique por qué se eligió en función de los datos
% y del mensaje que se quiere transmitir.

[Ejemplo: ``Se utilizó una escala secuencial de rojos (de claro a oscuro) para
representar la intensidad creciente de la variable. El rojo refuerza la
connotación de riesgo o urgencia asociada al macrotema, y la escala secuencial
es adecuada para datos continuos sin un punto medio significativo...'']

\subsection{Tipografía empleada}

% Indique la fuente usada en el mapa (títulos, etiquetas, leyenda) y cómo aporta
% coherencia al diseño general de la infografía.

[Ejemplo: ``Se utilizó \textit{Inter} para títulos y etiquetas del mapa, coherente
con la tipografía de la infografía grupal. Su diseño sans-serif garantiza
legibilidad a distintas escalas y no compite visualmente con los elementos
geográficos del mapa.'']

\subsection{Mapa}

\begin{figure}[H]
  \centering
  \includegraphics[width=\textwidth]{figures/figure001.png}
  \caption{[Tipo de mapa]: [descripción del contenido]. Fuente: [fuente del dataset].}
  \label{fig:mapa1}
\end{figure}

\subsection{Conclusiones}

\begin{itemize}
  \item [Conclusión 1: patrón o hallazgo principal visible en el mapa.]
  \item [Conclusión 2: comparación geográfica relevante o anomalía destacable.]
  \item [Conclusión 3: relación con el macrotema o implicancia para el público objetivo.]
\end{itemize}

% ─────────────────────────────────────────────────────────────────────────────
\section{Trabajo individual -- Integrante 2: [Nombre]}

\subsection{Dataset seleccionado}

\begin{itemize}
  \item \textbf{Nombre del dataset:} [nombre]
  \item \textbf{Fuente:} [URL o referencia]
  \item \textbf{Variables relevantes:} [lista breve]
  \item \textbf{Cobertura:} [países, período temporal, granularidad geográfica]
  \item \textbf{Descripción cualitativa:} [2--3 oraciones describiendo el dataset]
\end{itemize}

\subsection{Objetivo de la visualización}

\begin{itemize}
  \item \textbf{Pregunta que busca responder:} [¿Qué pregunta responde el mapa?]
  \item \textbf{Público objetivo:} [¿A quién está dirigido?]
  \item \textbf{Acción o decisión que podría apoyar:} [¿Qué decisión informaría?]
\end{itemize}

\subsection{Tipo de mapa y herramienta}

\begin{itemize}
  \item \textbf{Tipo de mapa:} [Coropleta / Puntos / Burbujas / Flujo / Calor / otro]
  \item \textbf{Herramienta:} [Flourish / Carto / Kepler.gl / Datawrapper / otra]
  \item \textbf{Justificación del tipo:} [¿Por qué este tipo de mapa es el más adecuado?]
\end{itemize}

\subsection{Mapa}

\begin{figure}[H]
  \centering
  \includegraphics[width=\textwidth]{figures/figure001.png}
  \caption{[Tipo de mapa]: [descripción del contenido]. Fuente: [fuente del dataset].}
  \label{fig:mapa2}
\end{figure}

\subsection{Conclusiones}

\begin{itemize}
  \item [Conclusión 1.]
  \item [Conclusión 2.]
  \item [Conclusión 3.]
\end{itemize}

% ─────────────────────────────────────────────────────────────────────────────
\section{Conclusiones generales}

[Párrafo integrando los hallazgos de los tres mapas en relación al macrotema.
¿Qué historia cuentan los tres mapas en conjunto? ¿Qué patrones geográficos
emergen? ¿Qué preguntas quedan abiertas para una próxima entrega?]

% ─────────────────────────────────────────────────────────────────────────────
\section{Repositorio de GitHub}

El código fuente y los archivos de esta entrega se encuentran disponibles en:

\begin{center}
  \href{https://github.com/[usuario]/[repositorio]}{\texttt{https://github.com/[usuario]/[repositorio]}}
\end{center}

\begin{verbatim}
/
|-- Informe_3
|   |-- figures/
|   |   |-- infografia.png
|   |   |-- mapa1.png
|   |   `-- mapa2.png
|   |-- informe.pdf
|   `-- informe.tex
`-- README.md
\end{verbatim}

% ─────────────────────────────────────────────────────────────────────────────

\end{document}